\chapter{Cross-Cloud Service Equivalence}
\index{cross-cloud equivalence}\index{service mapping}\index{multi-cloud}%
This appendix is a living reference that maps conceptually equivalent services and
ideas across the seven platforms studied in this book series: Amazon Web Services
(AWS), Microsoft Azure, Databricks, Microsoft Fabric, Google Cloud Platform (GCP),
MongoDB, and Snowflake. The names differ, but the underlying data-engineering
concepts---durable object storage, tiered lifecycle economics, immutability, managed
relational engines, elastic compute, and disciplined data organization---recur
everywhere.

\begin{notebox}
This appendix grows with each chapter. Every new chapter contributes the equivalence
tables introduced in its cross-cloud callout, so by the end of the book you will hold
a single consolidated Rosetta Stone for moving fluently between clouds. Where a clean
one-to-one mapping does not exist, the prose notes the closest analog and the caveat.
\end{notebox}

\section{Object and Data-Lake Storage}
\index{object storage!equivalence}\index{storage classes!equivalence}%
The foundational layer covered in Chapter~1. Every platform offers durable, HTTP-addressable
object storage with tiered pricing, lifecycle transitions, versioning, and write-once-read-many
(WORM) controls.

\begin{center}
\footnotesize
\begin{tabular}{@{}p{2.8cm}p{3.2cm}p{3.2cm}p{3.2cm}@{}}
\toprule
\textbf{Capability} & \textbf{AWS} & \textbf{Azure} & \textbf{GCP} \\
\midrule
Object storage & Amazon S3 & ADLS Gen2 / Blob & Cloud Storage \\
Hot tier & S3 Standard & Hot & Standard \\
Infrequent access & S3 Standard-IA & Cool & Nearline \\
Archive (instant) & Glacier Instant Retrieval & Cold & Coldline \\
Deep archive & Glacier Deep Archive & Archive & Archive \\
Lifecycle rules & Lifecycle policies & Lifecycle management & Object Lifecycle Mgmt \\
Versioning & S3 Versioning & Blob versioning & Object Versioning \\
Immutability (WORM) & Object Lock & Immutable blob storage & Bucket Lock \\
Cross-region copy & CRR / SRR & Object replication & Dual/multi-region + Transfer \\
Encryption at rest & SSE-S3 / SSE-KMS & SSE + CMK (Key Vault) & Google-managed / CMEK \\
\bottomrule
\end{tabular}
\end{center}

\noindent The three remaining platforms approach storage differently.
\textbf{Microsoft Fabric} exposes a single logical lake, \emph{OneLake}, and uses
\emph{shortcuts} to virtualize S3, ADLS Gen2, and Google Cloud Storage in place without copying.
\textbf{Snowflake} abstracts storage entirely: table data lives in cloud object storage as
immutable, columnar \emph{micro-partitions} that the service manages on your behalf.
\textbf{MongoDB} persists documents through the \emph{WiredTiger} storage engine, balancing an
internal cache against the operating-system page cache rather than exposing storage tiers.

\section{Managed Relational Databases}
\index{relational databases!equivalence}%
The focus of Chapter~2 for the three hyperscalers. Each provides a fully managed relational
engine, a cloud-native high-performance variant, high availability, read replicas, and a
migration service.

\begin{center}
\footnotesize
\begin{tabular}{@{}p{2.8cm}p{3.2cm}p{3.2cm}p{3.2cm}@{}}
\toprule
\textbf{Capability} & \textbf{AWS} & \textbf{Azure} & \textbf{GCP} \\
\midrule
Managed relational & Amazon RDS & Azure SQL Database & Cloud SQL \\
Cloud-native engine & Amazon Aurora & SQL Hyperscale & AlloyDB \\
Managed PostgreSQL & RDS / Aurora PostgreSQL & Azure DB for PostgreSQL & Cloud SQL / AlloyDB \\
High availability & Multi-AZ & Zone-redundant HA & Regional HA \\
Read scaling & Read Replicas & Read replicas / geo-replicas & Read Replicas \\
Global / geo & Aurora Global Database & Active geo-replication & Cross-region replicas \\
Serverless compute & Aurora Serverless v2 & Azure SQL Serverless & AlloyDB (elastic) \\
Migration service & DMS + SCT & Azure Database Migration Svc & Database Migration Svc \\
\bottomrule
\end{tabular}
\end{center}

\noindent \textbf{Snowflake}, \textbf{Fabric}, and \textbf{MongoDB} are not managed relational
engines, but they solve overlapping problems: Snowflake and Fabric provide SQL analytics over
warehouse/lakehouse storage, while MongoDB provides a distributed document store whose schema
design (Chapter~2) replaces rigid relational normalization with intentional, access-driven modeling.

\section{Elastic Compute and Analytical Warehouses}
\index{elastic compute!equivalence}\index{virtual warehouse!equivalence}%
Snowflake's Chapter~2 topic---separating compute from storage and scaling it elastically---has a
counterpart in every analytics platform.

\begin{center}
\footnotesize
\begin{tabular}{@{}p{2.8cm}p{3.2cm}p{3.2cm}p{3.2cm}@{}}
\toprule
\textbf{Concept} & \textbf{Snowflake} & \textbf{AWS} & \textbf{Azure / GCP} \\
\midrule
Compute unit & Virtual Warehouse & Redshift RA3 / Serverless & Synapse pool / BigQuery slots \\
Elastic concurrency & Multi-cluster warehouse & Concurrency scaling & Autoscale / on-demand \\
Pause when idle & Auto-suspend / resume & Pause / resume & Auto-pause / serverless \\
Billing unit & Credits & RPU / node-hours & DWU / slot-hours \\
Result reuse & Result cache & Result cache & Result / BI cache \\
\bottomrule
\end{tabular}
\end{center}

\noindent \textbf{Microsoft Fabric} meters all workloads against a shared pool of Capacity Units
(CU) on an F-SKU, an approach closer to a single elastic budget than to per-engine warehouses.
\textbf{MongoDB Atlas} scales through cluster tiers with optional auto-scaling of compute and storage.

\section{Data Organization and the Medallion Pattern}
\index{medallion architecture!equivalence}%
Fabric's Chapter~2 topic---the Bronze/Silver/Gold medallion---is a portable design pattern rather
than a product. Every platform implements the same progressive-refinement idea with its own primitives.

\begin{center}
\footnotesize
\begin{tabular}{@{}p{2.6cm}p{3.4cm}p{3.4cm}p{3.0cm}@{}}
\toprule
\textbf{Layer} & \textbf{Fabric} & \textbf{Data lake (AWS/Azure/GCP)} & \textbf{Snowflake} \\
\midrule
Bronze (raw) & Bronze Lakehouse & \texttt{raw/} prefix (S3/ADLS/GCS) & RAW schema / stage \\
Silver (validated) & Silver Lakehouse & \texttt{validated/} prefix & STAGING schema \\
Gold (curated) & Gold Lakehouse & \texttt{curated/} prefix & ANALYTICS schema \\
Storage format & Delta-Parquet & Parquet / Delta / Iceberg & Micro-partitions \\
\bottomrule
\end{tabular}
\end{center}

\noindent \textbf{MongoDB} realizes the same separation with distinct collections or databases
(for example \texttt{raw\_events}, \texttt{validated}, \texttt{curated}) governed by schema-design
patterns rather than a lake layout.

\section{Document Data Modeling}
\index{document modeling!equivalence}%
MongoDB's Chapter~2 schema-design patterns generalize to every document and wide-column store.

\begin{center}
\footnotesize
\begin{tabular}{@{}p{3.0cm}p{3.2cm}p{3.2cm}p{3.0cm}@{}}
\toprule
\textbf{Concept} & \textbf{MongoDB} & \textbf{Amazon DynamoDB} & \textbf{Azure Cosmos DB} \\
\midrule
Primary access unit & Document / collection & Item / table & Item / container \\
Partitioning & Shard key & Partition key & Partition key \\
One-table modeling & Embedding + patterns & Single-table design & Embedding + partitioning \\
Time-series grouping & Bucket Pattern & Composite sort keys & Bucketing by key \\
Optional attributes & Attribute Pattern & Sparse attributes & Sparse properties \\
\bottomrule
\end{tabular}
\end{center}

\noindent Google Cloud's closest analogs are \textbf{Firestore} and \textbf{Bigtable}. In
\textbf{Snowflake} and \textbf{Fabric}, semi-structured documents are stored in \texttt{VARIANT}
(Snowflake) or JSON columns and queried with path expressions rather than modeled as collections.

\section{Globally Distributed and NoSQL Databases}
\index{NoSQL!equivalence}\index{distributed databases!equivalence}%
Chapter~3 branches by platform into high-scale NoSQL (AWS DynamoDB, Azure Cosmos DB), globally
consistent NewSQL (GCP Cloud Spanner), and advanced document modeling (MongoDB). They all answer
the same question: how to scale a database horizontally across partitions and regions.

\begin{center}
\footnotesize
\begin{tabular}{@{}p{2.8cm}p{2.9cm}p{2.9cm}p{2.9cm}p{2.4cm}@{}}
\toprule
\textbf{Concept} & \textbf{DynamoDB} & \textbf{Cosmos DB} & \textbf{Cloud Spanner} & \textbf{MongoDB} \\
\midrule
Model & Key-value / doc & Multi-model & Distributed SQL & Document \\
Partitioning & Partition key & Partition key & Split by key range & Shard key \\
Throughput unit & RCU / WCU & Request Units (RU) & Nodes / processing units & Cluster tier \\
Secondary indexes & GSI / LSI & Automatic indexing & Secondary indexes & Indexes \\
Change stream & DynamoDB Streams & Change Feed & Change streams & Change streams \\
Expiry & TTL & TTL & (manual) & TTL index \\
Multi-region & Global Tables & Multi-region writes & Multi-region (synchronous) & Zones / global clusters \\
Consistency & Eventual / strong & 5 tunable levels & External (TrueTime) & Tunable read/write concern \\
\bottomrule
\end{tabular}
\end{center}

\noindent \textbf{Snowflake} and \textbf{Fabric} are analytical platforms rather than OLTP key-value
stores; they ingest from these operational databases and store semi-structured payloads in
\texttt{VARIANT}/JSON columns. Google's dedicated NoSQL analogs are \textbf{Firestore} (document)
and \textbf{Bigtable} (wide-column).

\section{Managed Apache Spark}
\index{Apache Spark!equivalence}%
Microsoft Fabric's Chapter~3 topic---managed Spark pools and notebooks---has a first-class
equivalent on every cloud.

\begin{center}
\footnotesize
\begin{tabular}{@{}p{2.8cm}p{3.2cm}p{3.2cm}p{3.2cm}@{}}
\toprule
\textbf{Capability} & \textbf{Fabric} & \textbf{AWS} & \textbf{Azure / GCP} \\
\midrule
Managed Spark & Starter / Custom pools & Glue / EMR / EMR Serverless & Synapse Spark / Databricks; Dataproc \\
Notebook UX & Fabric Notebooks & Glue Studio / EMR Studio & Synapse / Databricks; Vertex AI \\
Serverless option & Autoscaling pools & Glue / EMR Serverless & Synapse serverless; Dataproc Serverless \\
Lake format & Delta on OneLake & Delta / Iceberg / Hudi & Delta; Delta / Iceberg \\
\bottomrule
\end{tabular}
\end{center}

\noindent \textbf{Snowflake} offers \textbf{Snowpark} for DataFrame-style processing, and
\textbf{MongoDB} integrates through the Spark Connector.

\section{Query Acceleration and Materialized Views}
\index{materialized views!equivalence}\index{query acceleration!equivalence}%
Snowflake's Chapter~3 acceleration tools---materialized views, the Search Optimization Service, and
transient/temporary tables---map to comparable features in every warehouse.

\begin{center}
\footnotesize
\begin{tabular}{@{}p{2.8cm}p{3.0cm}p{3.0cm}p{3.4cm}@{}}
\toprule
\textbf{Feature} & \textbf{Snowflake} & \textbf{AWS Redshift} & \textbf{Azure Synapse / GCP BigQuery} \\
\midrule
Materialized views & Materialized Views & Materialized views & Materialized views \\
Point-lookup accel. & Search Optimization Svc & Sort keys / zone maps & (indexing) / search indexes \\
Result reuse & Result cache & Result cache & Result-set cache / BI Engine \\
Ephemeral tables & Transient / temporary & Temp tables & Temp tables \\
\bottomrule
\end{tabular}
\end{center}

\noindent \textbf{Fabric} accelerates reads through Direct Lake and its Warehouse engine, while
\textbf{MongoDB} relies on secondary and compound indexes plus the in-memory working set.

\section{In-Memory Caching and Hybrid Ingestion}
\index{caching!equivalence}\index{in-memory cache!equivalence}%
The AWS and Azure Chapter~4 topic---low-latency caching in front of a durable store, fed by hybrid
(edge-to-cloud) ingestion---has a managed Redis-compatible equivalent on every cloud.

\begin{center}
\footnotesize
\begin{tabular}{@{}p{2.8cm}p{3.0cm}p{3.0cm}p{3.4cm}@{}}
\toprule
\textbf{Capability} & \textbf{AWS} & \textbf{Azure} & \textbf{GCP / others} \\
\midrule
Managed cache & ElastiCache (Redis/Memcached) & Cache for Redis & Memorystore (Redis/Memcached) \\
Real-time API layer & AppSync / API Gateway & API Management / SignalR & API Gateway / Apigee \\
Hybrid ingestion & DataSync / Snow family / IoT Core & Data Box / IoT Hub / Arc & Transfer Appliance / IoT Core \\
Streaming intake & Kinesis / MSK & Event Hubs & Pub/Sub \\
\bottomrule
\end{tabular}
\end{center}

\noindent \textbf{Snowflake} and \textbf{Fabric} rely on result and Direct~Lake caches rather than a
standalone Redis tier, while \textbf{MongoDB} serves hot data from its in-memory working set and
Atlas can pair with an external cache.

\section{Wide-Column and Key-Value NoSQL}
\index{wide-column!equivalence}\index{Bigtable!equivalence}%
GCP's Chapter~4 topic---Cloud Bigtable---is a wide-column, high-throughput, low-latency NoSQL store.
Its peers span the managed and open-source ecosystems.

\begin{center}
\footnotesize
\begin{tabular}{@{}p{2.8cm}p{3.0cm}p{3.0cm}p{3.4cm}@{}}
\toprule
\textbf{Characteristic} & \textbf{GCP Bigtable} & \textbf{AWS} & \textbf{Azure / open source} \\
\midrule
Data model & Wide-column & DynamoDB (key-value/doc) & Cosmos DB (multi-model) \\
HBase-compatible & HBase API & Keyspaces (Cassandra) & Managed Instance for Cassandra \\
Scaling unit & Nodes / clusters & On-demand / provisioned & RU/s (Cosmos) \\
Best fit & Time-series, IoT, ad-tech & Serverless key-value & Global multi-model \\
\bottomrule
\end{tabular}
\end{center}

\noindent \textbf{Snowflake} and \textbf{Fabric} address these workloads analytically rather than as
operational key-value stores; \textbf{MongoDB} covers the document end of the same NoSQL spectrum.

\section{Data Protection: Cloning, Time Travel, and Migration}
\index{cloning!equivalence}\index{time travel!equivalence}\index{point-in-time recovery!equivalence}%
Snowflake's Chapter~4 topic---zero-copy cloning, Time Travel, and Fail-safe---plus MongoDB's
relational-to-document migration map to data-protection and migration tooling across the clouds.

\begin{center}
\footnotesize
\begin{tabular}{@{}p{2.8cm}p{3.0cm}p{3.0cm}p{3.4cm}@{}}
\toprule
\textbf{Capability} & \textbf{Snowflake} & \textbf{AWS / Azure} & \textbf{GCP / Fabric} \\
\midrule
Instant clone & Zero-copy clone & Aurora clone / SQL DB copy & BigQuery table clone; OneLake shortcut \\
Point-in-time recovery & Time Travel & PITR (RDS/SQL DB backups) & PITR; Delta time travel \\
Retention safety net & Fail-safe (7 days) & Automated backups & Backups; Delta VACUUM window \\
Schema migration & --- & DMS / SCT & Database Migration Service \\
\bottomrule
\end{tabular}
\end{center}

\noindent \textbf{MongoDB} migrations use \textbf{Relational Migrator} plus the embed-vs-reference
patterns from Chapter~4, and \textbf{Delta Lake} exposes historical versions through
\texttt{VERSION AS OF} time-travel queries.

\section{Account Setup and Free Tiers}
\index{free tier!equivalence}\index{account setup!equivalence}%
Chapter~0's setup decisions---identity, CLI, and cost guardrails---have the same shape on every
cloud before a single byte is stored.

\begin{center}
\footnotesize
\begin{tabular}{@{}p{2.9cm}p{3.0cm}p{3.0cm}p{3.0cm}@{}}
\toprule
\textbf{Setup step} & \textbf{AWS} & \textbf{Azure} & \textbf{GCP} \\
\midrule
Free offer & Free Tier (12 mo + always free) & Free account (\$200/30 days + always free) & Free trial (\$300/90 days + Always Free) \\
Sign-up guard & Card + root email & Card + Microsoft account & Card + Google account \\
Identity model & IAM / IAM Identity Center & Microsoft Entra ID + RBAC & Cloud IAM \\
Admin isolation & Never use root daily & Avoid Global Admin daily & Avoid Owner; least privilege \\
CLI & AWS CLI v2 (\texttt{aws configure}) & Azure CLI (\texttt{az login}) & gcloud CLI (\texttt{gcloud init}) \\
Browser shell & AWS CloudShell & Azure Cloud Shell & Cloud Shell \\
Cost guardrail & AWS Budgets + alerts & Cost Management budgets & Budgets \& alerts \\
Console & AWS Management Console & Azure Portal & Google Cloud Console \\
\bottomrule
\end{tabular}
\end{center}

\noindent This book's platform, \textbf{MongoDB}, is used two ways---a forever-free Atlas M0
cluster (no card) secured by a database user and an IP access list, and a local Community Server
for storage-engine labs. \textbf{Microsoft Fabric} uses a Microsoft 365/Power BI work account plus
a trial capacity (no card); \textbf{Snowflake} uses a 30-day/\$400 trial (no card) guarded by
Resource Monitors; \textbf{Databricks} offers a free trial and a forever-free Community/Express
edition, with serverless compute billed to the cloud account or Databricks depending on the model.

\section{Consensus, Replication, and Durable Writes}
\index{consensus!equivalence}\index{replication!equivalence}%
Chapter~5's topic. Every managed database hides consensus behind a product name, but the
trade-offs---quorum size, failover time, rollback risk---are identical.

\begin{center}
\footnotesize
\begin{tabular}{@{}p{2.9cm}p{3.0cm}p{3.0cm}p{3.0cm}@{}}
\toprule
\textbf{Capability} & \textbf{AWS} & \textbf{Azure} & \textbf{GCP} \\
\midrule
Managed replication & RDS Multi-AZ & Azure SQL auto-failover groups & Cloud SQL HA \\
Quorum/consensus & Aurora 6-copy quorum & Cosmos DB quorum & Spanner Paxos \\
Multi-region reads & Aurora Global / ElastiCache & Cosmos DB multi-region & Spanner multi-region \\
Failover control & Manual or auto promote & Priority-based failover & Automatic regional \\
Durable-write analog & \texttt{synchronous} commit & Strong consistency bound & Commit quorum \\
\addlinespace
Snowflake / Fabric / Databricks & \multicolumn{3}{l}{Consensus invisible; durability delegated to cloud object storage (S3/ADLS/GCS) with cross-region replication as a DR feature} \\
\bottomrule
\end{tabular}
\end{center}

\noindent \textbf{MongoDB} exposes what the warehouses hide: you choose the quorum that
acknowledges a write (write concern) and the consistency a read observes (read concern)
\cite{mongodbReplicationDocs,ongaro2014raft}. \textbf{Databricks} sits with the warehouses here:
its control plane is managed and the Delta transaction log on object storage is the durability
boundary, so there is no user-tunable quorum.

\section{Horizontal Partitioning and Shard Keys}
\index{sharding!equivalence}\index{partitioning!equivalence}%
Chapter~6's topic. Horizontal partitioning exists everywhere under different names; the shard-key
decision is the same decision as a DynamoDB partition key, a Cosmos DB partition key, or a Bigtable
row-key design.

\begin{center}
\footnotesize
\begin{tabular}{@{}p{2.9cm}p{3.0cm}p{3.0cm}p{3.0cm}@{}}
\toprule
\textbf{Concept} & \textbf{AWS} & \textbf{Azure} & \textbf{GCP} \\
\midrule
Partition unit & DynamoDB partition & Cosmos DB logical partition & Bigtable tablet \\
Router & Client SDK (DynamoDB) & Gateway + partitions & Front end + tablet map \\
Split metadata & Partition map & Partition map & Master tablet \\
Hot-key risk & Hot partition throttling & 20 GB logical cap, hot RU & Hot tablet \\
Even distribution tool & Hash of partition key & Hash of partition key & Salting / key design \\
Rebalancing & Automatic split & Physical partition split & Tablet splits \\
\addlinespace
Snowflake / Fabric & \multicolumn{3}{l}{Micro-partitions replace user-visible sharding; clustering keys influence pruning, not routing} \\
Databricks & \multicolumn{3}{l}{Liquid clustering replaces explicit partition keys; layout is physical optimization, never a routing decision} \\
\bottomrule
\end{tabular}
\end{center}

\noindent \textbf{MongoDB}'s difference is explicitness: you pick the shard key, you can refine it
live, and you can mis-pick it badly enough to create a hot shard \cite{mongodbShardingDocs}. The
lakehouse platforms deliberately removed that decision; MongoDB keeps it because operational
routing---not just analytical pruning---depends on it.

\section{Shard-Key Evolution and Hot-Key Mitigation}
\index{shard key!evolution}\index{hot key!equivalence}%
Chapter~7's topic. Shard-key optimization is the same discipline as partition-key tuning on the
managed key-value stores: the key decides distribution, and a wrong key surfaces as throttling on
one partition while the rest idle.

\begin{center}
\footnotesize
\begin{tabular}{@{}p{2.9cm}p{3.0cm}p{3.0cm}p{3.0cm}@{}}
\toprule
\textbf{Capability} & \textbf{AWS} & \textbf{Azure} & \textbf{GCP} \\
\midrule
Key change after creation & Not possible (rebuild table) & Not possible (new container) & Not possible (new table) \\
Hot-key mitigation & Adaptive capacity + DAX & RU redistribution, synthetic keys & Tablet splits + salting \\
Skew observability & Contributor Insights & Partition-key-metrics logs & Key Visualizer \\
Balancing & Automatic partition split & Physical partition split & Automatic tablet rebalance \\
Cross-partition txn & TransactWrite (bounded) & Not supported across partitions & Spanner synchronous commit \\
\addlinespace
Snowflake / Fabric & \multicolumn{3}{l}{Clustering keys guide micro-partition layout; reclustering runs in background with credit cost} \\
Databricks & \multicolumn{3}{l}{\texttt{ALTER TABLE ... CLUSTER BY} changes liquid-clustering keys online; \texttt{OPTIMIZE} rewrites layout in the background} \\
\bottomrule
\end{tabular}
\end{center}

\noindent \textbf{MongoDB} is uniquely forgiving among the operational stores:
\texttt{refineCollectionShardKey} and \texttt{reshardCollection} allow online key evolution that
DynamoDB and Cosmos DB simply do not offer \cite{mongodbShardKeyDocs}. Databricks matches that
forgiveness on the analytical side by making layout rewritable without changing table identity.

\section{Data Residency and Global Distribution}
\index{data residency!equivalence}\index{global clusters!equivalence}%
Chapter~8's topic. Data residency is a compliance requirement that every cloud implements as a
routing primitive: data pinned to regions by policy.

\begin{center}
\footnotesize
\begin{tabular}{@{}p{2.9cm}p{3.0cm}p{3.0cm}p{3.0cm}@{}}
\toprule
\textbf{Capability} & \textbf{AWS} & \textbf{Azure} & \textbf{GCP} \\
\midrule
Residency primitive & DynamoDB global tables + Regions & Cosmos DB multi-region + policies & Spanner placements / Firestore multi-region \\
Local writes & Multi-region active-active & Multi-region writes & Multi-region configs \\
Read locality & Global tables nearest region & Cosmos DB preferred regions & Spanner read-only replicas \\
Residency enforcement & Region-scoped IAM + SCPs & Azure Policy + regions & Org policy constraints \\
Conflict handling & Last-writer-wins & LWW / custom merge & Synchronous commit \\
\addlinespace
Snowflake / Fabric & \multicolumn{3}{l}{Database replication between regions; residency by account placement rather than row-level pinning} \\
Databricks & \multicolumn{3}{l}{Workspace and Unity Catalog metastore pinned to a region; residency enforced at deployment boundary, not per row} \\
\bottomrule
\end{tabular}
\end{center}

\noindent \textbf{MongoDB}'s zone sharding pins \emph{ranges of data} to shards in named
regions---row-level residency with auditable chunk evidence---a stronger and more provable
guarantee than account-level regional placement \cite{mongodbZoneShardingDocs,mongodbGlobalClustersDocs}.

\section{Indexing and Physical Data Layout}
\index{indexing!equivalence}\index{physical layout!equivalence}%
Chapter~9's topic. Every analytical and operational store has an equivalent of ``index design
decides latency.''

\begin{center}
\footnotesize
\begin{tabular}{@{}p{2.9cm}p{3.0cm}p{3.0cm}p{3.0cm}@{}}
\toprule
\textbf{Capability} & \textbf{AWS} & \textbf{Azure} & \textbf{GCP} \\
\midrule
Managed document DB & DocumentDB & Cosmos DB (Mongo API) & Firestore / Bigtable \\
Index model & Manual indexes & Autopilot / manual policies & Composite indexes required \\
Secondary indexes & B-tree indexes & Automatic indexing (default) & Automatic single-field \\
Query plan insight & Profiler + explain & Diagnostic logs & Query explain stats \\
Columnar analytics & Redshift sort/dist keys & Dedicated SQL pool indexes & BigQuery clustering \\
Cloud data warehouse & Redshift distribution keys & Synapse distributions & BigQuery partitions \\
\addlinespace
Snowflake & \multicolumn{3}{l}{Micro-partition pruning, clustering keys, search optimization service} \\
Databricks & \multicolumn{3}{l}{Liquid clustering + Z-order, Photon vectorized execution, file-level data skipping on Delta statistics} \\
\bottomrule
\end{tabular}
\end{center}

\noindent The mental model transfers directly: Snowflake's pruning, BigQuery's clustering,
Redshift's sort keys, and Databricks' file skipping are all physical-data-layout decisions that
replace per-query index lookup with zone-map elimination. \textbf{MongoDB}'s compound B+Tree index
is the operational-system counterpart: created deliberately, charged to every write, and proven
with \texttt{explain()} rather than assumed \cite{mongodbIndexesDocs,mongodbExplainDocs}.

\section{Aggregation Pipelines and SQL Analytics}
\index{aggregation pipeline!equivalence}%
Chapter~10's topic. The aggregation pipeline is MongoDB's in-database analytics engine; its
siblings are the SQL engines of the cloud warehouses and lakehouses.

\begin{center}
\footnotesize
\begin{tabular}{@{}p{2.9cm}p{3.0cm}p{3.0cm}p{3.0cm}@{}}
\toprule
\textbf{Pipeline concept} & \textbf{AWS (Redshift)} & \textbf{Azure (Synapse/Fabric)} & \textbf{GCP (BigQuery)} \\
\midrule
\texttt{\$match} & \texttt{WHERE} & \texttt{WHERE} & \texttt{WHERE} \\
\texttt{\$group} & \texttt{GROUP BY} & \texttt{GROUP BY} & \texttt{GROUP BY} \\
\texttt{\$lookup} & \texttt{LEFT JOIN} & \texttt{LEFT JOIN} & \texttt{LEFT JOIN} \\
\texttt{\$facet} & Multiple CTEs in one scan & \texttt{UNION} of CTEs & Multi-branch subqueries \\
\texttt{\$graphLookup} & Recursive CTE & Recursive CTE & Recursive CTE / \texttt{WITH RECURSIVE} \\
\texttt{\$merge}/\texttt{\$out} & \texttt{INSERT INTO SELECT} & CTAS / \texttt{INSERT} & \texttt{CREATE TABLE AS} \\
\texttt{allowDiskUse} & Disk-based spill & Spill to tempdb & Spill to disk (slots) \\
\addlinespace
Snowflake & \multicolumn{3}{l}{One SQL engine covers all of the above; pruning replaces the early-\texttt{\$match} discipline automatically} \\
Databricks & \multicolumn{3}{l}{Spark SQL with the Catalyst optimizer and Photon; the optimizer, not the written order, owns the physical plan} \\
\bottomrule
\end{tabular}
\end{center}

\noindent The key difference is ownership: Snowflake, BigQuery, and Databricks optimize the
physical plan for you, while \textbf{MongoDB} executes stages largely in written order. Stage
order \emph{is} the physical plan, which makes pipeline literacy a first-class engineering skill
\cite{mongodbAggregationDocs}.

\section{Window Functions and Maintained Views}
\index{window functions!equivalence}\index{materialized views!maintenance}%
Chapter~11's topic. Window functions are the most portable analytical concept in this series:
SQL \texttt{OVER (PARTITION BY ... ORDER BY ...)} exists verbatim in Redshift, Synapse, BigQuery,
Snowflake, and Spark SQL. MongoDB's \texttt{\$setWindowFields} is the same mathematics with
document semantics.

\begin{center}
\footnotesize
\begin{tabular}{@{}p{2.9cm}p{3.0cm}p{3.0cm}p{3.0cm}@{}}
\toprule
\textbf{Concept} & \textbf{AWS (Redshift)} & \textbf{Azure (Synapse)} & \textbf{GCP (BigQuery)} \\
\midrule
Window stage & \texttt{OVER()} clause & \texttt{OVER()} clause & \texttt{OVER()} clause \\
Rank & \texttt{RANK}/\texttt{DENSE\_RANK} & Same & Same \\
Moving average & \texttt{AVG() OVER (ROWS)} & Same & Same \\
Derivative & \texttt{LAG()} arithmetic & Same & Same \\
Materialized view & Redshift MVs (auto-refresh) & Indexed views / Spark SQL & BigQuery MVs (auto-refresh) \\
Incremental refresh & Automatic & Manual / pipeline & Automatic, partition-aware \\
\addlinespace
Snowflake & \multicolumn{3}{l}{Window functions identical; materialized views auto-maintained; streams + tasks for incremental pipelines} \\
Databricks & \multicolumn{3}{l}{Identical SQL windows; Delta Live Tables materialized views refresh incrementally on a declared schedule or trigger} \\
\bottomrule
\end{tabular}
\end{center}

\noindent \textbf{MongoDB}'s distinction is the manual refresh path: \texttt{\$merge} into a
results collection is a materialized view you maintain yourself, so you own its freshness,
idempotency, and backfill story \cite{mongodbSetWindowFieldsDocs,mongodbMergeDocs}.

\section{Performance Diagnosis Instruments}
\index{query plan!equivalence}\index{slow log!equivalence}%
Chapter~12's topic. Every platform exposes the same four instruments---a query plan, a slow log,
wait/queue metrics, and per-object usage statistics---with different names.

\begin{center}
\footnotesize
\begin{tabular}{@{}p{2.9cm}p{3.0cm}p{3.0cm}p{3.0cm}@{}}
\toprule
\textbf{Instrument} & \textbf{AWS (RDS/Redshift)} & \textbf{Azure (SQL/Cosmos)} & \textbf{GCP (SQL/BigQuery)} \\
\midrule
Query plan & EXPLAIN / Redshift plan & Query Store plans & EXPLAIN / dry-run slots \\
Slow log & RDS slow query log & Azure SQL diagnostics & Cloud SQL insights \\
Wait/queue metrics & Performance Insights waits & DMVs / wait stats & Query Insights \\
Object usage stats & pg\_stat\_user\_indexes & sys.dm\_db\_index\_usage\_stats & INFORMATION\_SCHEMA \\
\addlinespace
Snowflake / Fabric & \multicolumn{3}{l}{\texttt{QUERY\_HISTORY}, query profile operator timings, warehouse load metrics replace per-index statistics} \\
Databricks & \multicolumn{3}{l}{Spark UI stages/tasks, query profile, and \texttt{system.query.history} / system tables for fleet-wide usage} \\
\bottomrule
\end{tabular}
\end{center}

\noindent \textbf{MongoDB}'s distinctive contribution is that all four instruments are queryable
data inside the database itself---\texttt{system.profile}, \texttt{explain()},
\texttt{serverStatus}, \texttt{\$indexStats}---so diagnosis is a scripting exercise, not a console
scavenger hunt \cite{mongodbProfilerDocs,mongodbExplainDocs}.

\section{Full-Text Search Engines}
\index{full-text search!equivalence}\index{Lucene!equivalence}%
Chapter~13's topic. Full-text search as a managed engine beside the operational store is a
universal pattern: each cloud pairs its database with a Lucene-descended search service.

\begin{center}
\footnotesize
\begin{tabular}{@{}p{2.9cm}p{3.0cm}p{3.0cm}p{3.0cm}@{}}
\toprule
\textbf{Capability} & \textbf{AWS} & \textbf{Azure} & \textbf{GCP} \\
\midrule
Search service & OpenSearch & Azure AI Search & Vertex AI Search \\
Engine lineage & Lucene & Lucene & Proprietary + vector \\
Sync mechanism & DMS / streams / Lambda & Indexers / push API & Data connectors \\
Analyzers/tokenizers & OpenSearch analyzers & Lucene + custom analyzers & Managed, less exposed \\
Faceting & Aggregations & Facets & Built-in \\
Typo tolerance & Fuzzy queries & Fuzzy + synonyms & Managed \\
\addlinespace
Snowflake / Fabric & \multicolumn{3}{l}{Snowflake Search Optimization + Cortex Search (managed hybrid); Fabric search via OneLake + AI Search} \\
Databricks & \multicolumn{3}{l}{No standalone text engine; pairs Mosaic AI Vector Search with keyword pre-filtering, or integrates Azure AI Search / OpenSearch} \\
\bottomrule
\end{tabular}
\end{center}

\noindent Atlas Search's distinctive property is integration depth: the Lucene index lives inside
the Atlas platform, syncs from the collection automatically via change streams, and is queried
through the same aggregation pipeline as everything else \cite{mongodbAtlasSearchDocs,apacheLuceneRepo}.

\section{Vector Search and Retrieval-Augmented Generation}
\index{vector search!equivalence}\index{RAG!equivalence}%
Chapter~14's topic. Vector search moved from specialty to commodity in two years; the
differentiators are index algorithm, pre-filtering, and operational integration with the primary
store.

\begin{center}
\footnotesize
\begin{tabular}{@{}p{2.9cm}p{3.0cm}p{3.0cm}p{3.0cm}@{}}
\toprule
\textbf{Capability} & \textbf{AWS} & \textbf{Azure} & \textbf{GCP} \\
\midrule
Vector store & OpenSearch k-NN / pgvector on Aurora & Cosmos DB vector search / AI Search & Vertex AI Vector Search / AlloyDB \\
Embedding service & Bedrock Titan & Azure OpenAI & Vertex AI embeddings \\
Hybrid search & OpenSearch BM25 + k-NN & AI Search hybrid & Vertex hybrid ranking \\
RAG framework & Bedrock Knowledge Bases & Azure OpenAI + AI Search & Vertex AI RAG Engine \\
Co-located OLTP & Aurora pgvector & Cosmos DB (same doc) & AlloyDB (same row) \\
\addlinespace
Snowflake / Fabric & \multicolumn{3}{l}{Snowflake Cortex (embed + vector similarity in SQL); Fabric integrates with AI Search via OneLake} \\
Databricks & \multicolumn{3}{l}{Mosaic AI Vector Search over Delta tables + Foundation Model APIs; governance through Unity Catalog} \\
\bottomrule
\end{tabular}
\end{center}

\noindent Atlas Vector Search's distinctive property is co-location: vectors, BM25 text indexes,
and the operational documents live in one system, so hybrid retrieval and its filters execute in
one aggregation pipeline without a second store to sync
\cite{mongodbAtlasVectorSearch,mongodbAtlasSearchDocs,malkov2020hnsw}.

\section{Stream Processing}
\index{stream processing!equivalence}%
Chapter~15's topic. Managed streaming with SQL-like or pipeline semantics over a broker backbone
is a commodity layer; the differentiators are state management, exactly-once composition, and how
tightly the sink integrates with the operational store.

\begin{center}
\footnotesize
\begin{tabular}{@{}p{2.9cm}p{3.0cm}p{3.0cm}p{3.0cm}@{}}
\toprule
\textbf{Capability} & \textbf{AWS} & \textbf{Azure} & \textbf{GCP} \\
\midrule
Broker backbone & MSK / Kinesis & Event Hubs (Kafka API) & Pub/Sub / Managed Kafka \\
Stream processor & Kinesis Data Analytics (Flink) & Stream Analytics / Fabric RTI & Dataflow (Beam) \\
Windowing & Tumbling/hopping/session & Tumbling/hopping/session & Full Beam window model \\
State/checkpointing & Flink state backends & Managed & Managed (streaming engine) \\
Exactly-once sink & Idempotent/transactional sinks & Idempotent upserts & Idempotent writes \\
DB-native option & DMS streams & Cosmos change feed + Functions & Firestore triggers + Dataflow \\
\addlinespace
Snowflake / Fabric & \multicolumn{3}{l}{Streams + Tasks for incremental CDC pipelines; Snowpipe Streaming; Fabric eventstreams with KQL windows} \\
Databricks & \multicolumn{3}{l}{Structured Streaming (micro-batch) with checkpointed state and \texttt{MERGE} sinks into Delta; \texttt{availableNow} unifies batch and stream} \\
\bottomrule
\end{tabular}
\end{center}

\noindent Atlas Stream Processing's distinction is that the streaming tier speaks the aggregation
pipeline and sinks directly into Atlas collections, so enrichment joins and materialized outputs
stay on one platform \cite{mongodbStreamProcessingDocs,apacheKafkaRepo}.

\section{Change Data Capture, Triggers, and Reverse ETL}
\index{change data capture!equivalence}\index{reverse ETL!equivalence}%
Chapter~16's topic. CDC, serverless triggers, and reverse ETL are the connective tissue of every
modern data platform; the discipline---checkpoint, idempotency, replay---is identical everywhere.

\begin{center}
\footnotesize
\begin{tabular}{@{}p{2.9cm}p{3.0cm}p{3.0cm}p{3.0cm}@{}}
\toprule
\textbf{Concept} & \textbf{AWS} & \textbf{Azure} & \textbf{GCP} \\
\midrule
Database CDC & DynamoDB Streams / DMS CDC & Cosmos DB change feed & Firestore triggers / Datastream \\
Serverless trigger & Lambda on streams & Functions on change feed & Cloud Functions / Run \\
HTTPS data API & DynamoDB API / RDS Data API & Cosmos DB REST & Firestore REST \\
Reverse ETL & AppFlow / custom & Synapse Link + Logic Apps & Datastream + Workflows \\
Event backbone & Kinesis / EventBridge & Event Hubs / Event Grid & Pub/Sub \\
\addlinespace
Snowflake / Fabric & \multicolumn{3}{l}{Streams + Tasks for CDC; reverse ETL into operational stores via connectors; Fabric Activator for triggers} \\
Databricks & \multicolumn{3}{l}{Lakeflow Connect for managed CDC ingestion; \texttt{MERGE} / DLT \texttt{AUTO CDC} to apply changes into Delta} \\
\bottomrule
\end{tabular}
\end{center}

\noindent \textbf{MongoDB} change streams are distinguished by what they ride on: the same oplog
that powers replication, so CDC is a first-class, total-order-per-shard,
transactional-integrity-preserving API rather than a bolt-on log scrape
\cite{mongodbChangeStreamsDocs,mongodbOplogDocs}.

\section{Identity, Authentication, and Authorization}
\index{authentication!equivalence}\index{RBAC!equivalence}%
Chapter~17's topic. Identity and access design is the same discipline on every cloud: central
identity, least privilege, auditable grants.

\begin{center}
\footnotesize
\begin{tabular}{@{}p{2.9cm}p{3.0cm}p{3.0cm}p{3.0cm}@{}}
\toprule
\textbf{Capability} & \textbf{AWS} & \textbf{Azure} & \textbf{GCP} \\
\midrule
Cloud identity & IAM users/roles & Microsoft Entra ID & Cloud IAM \\
Workload identity & IAM roles for pods/EC2 & Managed identities & Workload Identity Federation \\
DB-level RBAC & DocumentDB users / IAM DB auth & Cosmos DB RBAC & Firestore security rules \\
Directory integration & AD Connector / LDAP & Native Entra & Cloud Identity \\
Audit & CloudTrail + engine audit & Diagnostic settings & Cloud Audit Logs \\
Certificate auth & ACM-issued client certs & Key Vault certs & Certificate Authority Service \\
\addlinespace
Snowflake / Fabric & \multicolumn{3}{l}{SCIM + SSO (SAML/OIDC), role hierarchy with inheritance, network policies; Fabric layers on Entra conditional access} \\
Databricks & \multicolumn{3}{l}{Account-console identities with SCIM/SSO; Unity Catalog grants (catalog/schema/table/column) as the data-plane RBAC} \\
\bottomrule
\end{tabular}
\end{center}

\noindent \textbf{MongoDB}'s distinctive surface is its \emph{in-database} RBAC: privileges named
by resource (database.collection) and action (down to pipeline-stage granularity), composable into
custom roles---closer to Snowflake's role model and Databricks' Unity Catalog grants than to IAM
policies \cite{mongodbSecurityChecklistDocs,mongodbRBACDocs}.

\section{The Encryption Ladder}
\index{encryption!equivalence}\index{key management!equivalence}%
Chapter~18's topic. Every cloud offers the same ladder---at-rest keys, TLS, client-side field
encryption---but the in-use layer is where vendors genuinely differ.

\begin{center}
\footnotesize
\begin{tabular}{@{}p{2.9cm}p{3.0cm}p{3.0cm}p{3.0cm}@{}}
\toprule
\textbf{Layer} & \textbf{AWS} & \textbf{Azure} & \textbf{GCP} \\
\midrule
Key management & KMS / CloudHSM & Key Vault / Managed HSM & Cloud KMS / HSM \\
At-rest & EBS/S3 SSE, RDS TDE & Disk + TDE & Default + CMEK \\
In transit & TLS everywhere & TLS + private endpoints & TLS + PSC \\
In use & SDK client-side encryption & Always Encrypted (SQL) & Confidential computing \\
Private connectivity & PrivateLink & Private Link & Private Service Connect \\
Erasure primitive & KMS key deletion & Key Vault purge & Key destruction \\
\addlinespace
Snowflake / Fabric & \multicolumn{3}{l}{Tri-Secret Secure (composite keys), column-level policies; customer-managed keys for the account layer} \\
Databricks & \multicolumn{3}{l}{Customer-managed keys for workspace storage and managed services; no field-level queryable encryption---pair with the operational store} \\
\bottomrule
\end{tabular}
\end{center}

\noindent \textbf{MongoDB}'s Queryable Encryption is the rare in-use implementation that keeps
\emph{equality queries} server-side without the server ever seeing plaintext---the cryptographic
boundary sits in the client driver, keyed by a KMS you control
\cite{mongodbQueryableEncryptionDocs,mongodbCSFLEDocs}.

\section{Private Networking and Audit Streaming}
\index{private endpoint!equivalence}\index{audit log!equivalence}%
Chapter~19's topic. Private connectivity and audit streaming are table stakes on every cloud; the
vocabulary differs, the architecture does not.

\begin{center}
\footnotesize
\begin{tabular}{@{}p{2.9cm}p{3.0cm}p{3.0cm}p{3.0cm}@{}}
\toprule
\textbf{Capability} & \textbf{AWS} & \textbf{Azure} & \textbf{GCP} \\
\midrule
Private endpoint & PrivateLink & Private Link & Private Service Connect \\
Coarse network gate & Security groups / NACLs & NSGs / firewall rules & VPC firewall rules \\
Peering & VPC peering / TGW & VNet peering & VPC Network Peering \\
TLS policy & ACM + TLS policies & Key Vault + TLS policy & Certificate Manager \\
Audit stream & CloudTrail + engine logs & Diagnostic settings + Entra logs & Cloud Audit Logs \\
SIEM targets & Security Hub / Splunk & Microsoft Sentinel & Chronicle / Splunk \\
\addlinespace
Snowflake / Fabric & \multicolumn{3}{l}{Network policies + private connectivity; access history views for audit; Fabric routes through Entra and Purview} \\
Databricks & \multicolumn{3}{l}{Private connectivity per cloud (PrivateLink/Private Link/PSC); audit logs delivered to storage and queryable via system tables} \\
\bottomrule
\end{tabular}
\end{center}

\noindent \textbf{MongoDB}'s specifics worth mastering: private endpoints into Atlas remove public
exposure by construction, and the auditing framework emits filterable JSON events---identity, DDL,
and scoped DML---that stream to your SIEM \cite{mongodbPrivateLinkDocs,mongodbAuditingDocs}.

\section{Backup, Tiering, and Queryable Archives}
\index{backup!equivalence}\index{archival!equivalence}%
Chapter~20's topic. Backup and tiering are solved problems at every vendor---the differentiator
is whether restores are drilled and whether the cold tier stays queryable.

\begin{center}
\footnotesize
\begin{tabular}{@{}p{2.9cm}p{3.0cm}p{3.0cm}p{3.0cm}@{}}
\toprule
\textbf{Capability} & \textbf{AWS} & \textbf{Azure} & \textbf{GCP} \\
\midrule
Managed backup & RDS snapshots + PITR & Automated backups + PITR & Cloud SQL PITR \\
Continuous/PITR & 35-day window & 7--35 days & 7--35 days \\
Cold tiering & S3 Glacier & Blob Archive & Archive Storage \\
Query over archive & Athena / Redshift Spectrum & Synapse / Fabric OneLake & BigQuery external tables \\
Long-term retention & Glacier Deep Archive & Archive tier & Archive class \\
Restore testing & Manual / scripted drills & Same & Same \\
\addlinespace
Snowflake / Fabric & \multicolumn{3}{l}{Time Travel + Fail-safe replace user-managed backups; Fabric OneLake shortcuts play the federated-query role} \\
Databricks & \multicolumn{3}{l}{Delta time travel (\texttt{VERSION AS OF}) and \texttt{RESTORE} for short windows; durability inherited from object storage, deep archive delegated to it} \\
\bottomrule
\end{tabular}
\end{center}

\noindent Atlas's version of the pattern: continuous cloud backups with PITR granularity inside
the oplog window, Online Archive tiering cold data to object storage as Parquet, and Atlas Data
Federation querying live and archived data in one statement
\cite{mongodbAtlasBackupDocs,mongodbOnlineArchiveDocs,mongodbDataFederationDocs}.

\section{Infrastructure as Code and Schema Change Management}
\index{infrastructure as code!equivalence}\index{Terraform!equivalence}%
Chapter~21's topic. Terraform manages all seven platforms, and the GitOps discipline---plan in
CI, apply on merge, state locked remotely---is identical everywhere.

\begin{center}
\footnotesize
\begin{tabular}{@{}p{2.9cm}p{3.0cm}p{3.0cm}p{3.0cm}@{}}
\toprule
\textbf{Capability} & \textbf{AWS} & \textbf{Azure} & \textbf{GCP} \\
\midrule
Native IaC & CloudFormation / CDK & ARM / Bicep & Deployment Manager \\
Terraform provider & \texttt{hashicorp/aws} & \texttt{azurerm} & \texttt{google} \\
DB migration tools & Flyway / Liquibase / Alembic & Same & Same \\
Plan-in-CI gate & \texttt{terraform plan} on PR & Same & Same \\
State backend & S3 + DynamoDB lock & Azure Blob + lease & GCS + lock \\
Secrets & Secrets Manager + CI vars & Key Vault + CI vars & Secret Manager + CI vars \\
\addlinespace
Snowflake / Fabric & \multicolumn{3}{l}{Snowflake provider for accounts/warehouses/grants; schemachange for migrations; Fabric via Terraform azapi and Git integration} \\
Databricks & \multicolumn{3}{l}{\texttt{databricks} provider for workspaces/catalogs/jobs; Asset Bundles package code + infrastructure as one versioned unit} \\
\bottomrule
\end{tabular}
\end{center}

\noindent \textbf{MongoDB}'s distinctive gap and opportunity: schema is not enforced by the
engine, so \emph{change management}---validators, indexes, migrations---is the mechanism that
keeps forty services' implicit schemas from drifting apart \cite{mongodbTerraformDocs,terraformDocs}.

\section{Kubernetes and Database Operators}
\index{Kubernetes!equivalence}\index{operator pattern!equivalence}%
Chapter~22's topic. Managed Kubernetes is commoditized across clouds; the
database-on-Kubernetes question---operator versus managed service---is asked identically
everywhere.

\begin{center}
\footnotesize
\begin{tabular}{@{}p{2.9cm}p{3.0cm}p{3.0cm}p{3.0cm}@{}}
\toprule
\textbf{Capability} & \textbf{AWS} & \textbf{Azure} & \textbf{GCP} \\
\midrule
Managed K8s & EKS & AKS & GKE \\
Operator hosting & EKS + operators & AKS + operators & GKE + operators (Autopilot caveats) \\
Stateful workload primitive & EBS-backed PVCs & Azure Disk PVCs & PD-backed PVCs \\
DBaaS alternative & DocumentDB / RDS & Cosmos DB & Atlas on GCP / Firestore \\
Secrets & Secrets Manager + ESO & Key Vault CSI/ESO & Secret Manager + ESO \\
\addlinespace
Snowflake / Fabric & \multicolumn{3}{l}{N/A as SaaS---but Snowpark Container Services runs containerized services next to data, the inverse direction} \\
Databricks & \multicolumn{3}{l}{N/A as SaaS---compute runs on the vendor-managed or customer-managed K8s underneath, invisible to the user by design} \\
\bottomrule
\end{tabular}
\end{center}

\noindent The operator pattern itself---declare desired state as a custom resource, let a
controller reconcile---is the industry default for stateful software on Kubernetes
\cite{mongodbK8sOperatorDocs,kubernetesDocs}. MongoDB is unusual in offering both poles: the
self-managed Community/Enterprise Operators and Atlas, which removes the cluster you would operate.

\section{Backup Services, Chaos, and Disaster Recovery}
\index{disaster recovery!equivalence}%
Chapter~23's topic. Every cloud sells regional redundancy, and every serious organization
discovers that recovery is a practiced behavior, not a purchased feature.

\begin{center}
\footnotesize
\begin{tabular}{@{}p{2.9cm}p{3.0cm}p{3.0cm}p{3.0cm}@{}}
\toprule
\textbf{Capability} & \textbf{AWS} & \textbf{Azure} & \textbf{GCP} \\
\midrule
Cross-region DB & Aurora Global / DDB global tables & Cosmos DB multi-region & Spanner multi-region \\
Backup service & AWS Backup & Azure Backup & Backup and DR Service \\
Chaos tooling & Fault Injection Service & Azure Chaos Studio & Third-party / Gremlin \\
Pilot-light DR & Warm standby patterns & Same & Same \\
Runbook automation & Systems Manager runbooks & Automation runbooks & Cloud Workflows \\
\addlinespace
Snowflake / Fabric & \multicolumn{3}{l}{Failover groups / database replication; RTO dominated by re-pointing clients and re-warming caches} \\
Databricks & \multicolumn{3}{l}{Workspace-level DR is a customer runbook (IaC rebuild + Delta deep clone); no managed cross-region workspace failover} \\
\bottomrule
\end{tabular}
\end{center}

\noindent \textbf{MongoDB}'s distinctive position: replica sets and sharding (Part~II) already
provide most of DR's mechanics---the chapter's work is measuring them under real failure and
writing the procedures humans execute when automation runs out \cite{mongodbOpsManagerDocs,sreBook}.

\section{The Whole Platform, End to End}
\index{reference architecture!equivalence}%
Chapter~24's synthesis. The final architecture is deliberately platform-shaped rather than
vendor-shaped: the same skeleton---declarative infrastructure, versioned migrations, CDC-fed
analytics, layered security, drilled recovery---rebuilds on any cloud with only control-plane
names changing.

\begin{center}
\footnotesize
\begin{tabular}{@{}p{2.9cm}p{3.0cm}p{3.0cm}p{3.0cm}@{}}
\toprule
\textbf{Platform layer} & \textbf{AWS flavor} & \textbf{Azure flavor} & \textbf{GCP flavor} \\
\midrule
Operational store & DocumentDB / Atlas & Cosmos DB / Atlas & Firestore / Atlas \\
IaC control plane & Terraform + AWS & Terraform + azurerm & Terraform + google \\
K8s staging & EKS + operator & AKS + operator & GKE + operator \\
CDC backbone & MSK / Kinesis & Event Hubs & Pub/Sub \\
Analytics sink & Redshift / S3+Athena & Fabric / Synapse & BigQuery \\
Governance & CloudTrail + Config & Purview + Policy & Dataplex + Audit Logs \\
\addlinespace
Databricks & \multicolumn{3}{l}{Plays the analytics-sink and AI layer on any row: Delta lakehouse + Unity Catalog + Mosaic AI, deployed on the same cloud's primitives} \\
\bottomrule
\end{tabular}
\end{center}

\noindent If you can defend the MongoDB version of this platform with evidence, you can rebuild it
on any row of that table---which is precisely the competence the certification and the interview
loop test for \cite{mongodbUniversityDocs,kleppmann2017ddia}.

\begin{tipbox}
Use this appendix as a study aid for the book's lockstep, cross-cloud approach: when you learn a
concept on one platform, immediately locate its equivalent here and predict how the other six
would express the same idea. That habit turns seven separate vocabularies into one mental model.
\end{tipbox}
